\subsection{Original product generators and genuine single-factor chart ratios}

\begin{lemma}[Every original single-factor ratio occurs in the actual product chart]
\label{algebra:actual-product-rees-chart-ratios}
\leanok
\uses{algebra:actual-rees-degree-zero-chart-equivalence}
Let $I$ be a finite index type and let $G_i:J_i\to A$ be arbitrary ORIGINAL
families over a commutative ring $A$.  Define the literal original product
family $P_J=\prod_iG_i(J_i)$.  For EVERY actual coordinate replacement,
\[
 P_{J[i:=k]}G_i(J_i)=G_i(k)P_J.
\]
This whole original numerator identity needs no domain or nonzero assumption.
If $A$ is a DOMAIN, nonzero factors give a nonzero product and a nonzero
product has every selected factor nonzero.  If $P_J\ne0$, then in the SAME
original fraction field
\[
 P_{J[i:=k]}/P_J=G_i(k)/G_i(J_i).
\]
Consequently every original single-factor ratio belongs to the ACTUAL
$J$-chart fraction subalgebra of the ORIGINAL product-family Rees scheme.
The individual index types $J_i$ need not be finite or nonempty; whenever a
selection $J$ is present, all statements retain its actual original elements.
\end{lemma}
\begin{proof}
Separate the $i$th factor from the finite product.  The dependent coordinate
replacement has the same original factors at every other index, so the two
displayed whole numerators are respectively
$G_i(k)(\prod_{l\ne i}G_l(J_l))G_i(J_i)$ and
$G_i(k)G_i(J_i)(\prod_{l\ne i}G_l(J_l))$; commutativity gives their equality
over any commutative ring, including rings with zero divisors.

Over the original domain, the actual finite-product nonzero criterion gives
both nonzero assertions.  Injectivity of the original coefficient map into
the SAME fraction field preserves nonzeroness of $P_J$ and $G_i(J_i)$.
The whole numerator identity can therefore be divided by the two actual
nonzero denominators, giving the exact original fraction equality.  Its left
side is an ORIGINAL coordinate ratio, hence a generator of the true product
chart's fraction-adjoin subalgebra.  Transport that membership through the
whole equality.  No principalization, center, chosen cover, radical or
fraction-comparison certificate is supplied.
\end{proof}
